\section{Morphological priming: Heyer and Kornishova's design and statistical detail}\label{sec:C}

This appendix carries the experimental detail from Heyer and Kornishova \autocite*{heyer-kornishova-2018-priming-eventually} that §\ref{sec:5} references in compressed form. The body keeps the design logic (short vs.~long SOA; transparency continuous; key result that transparency modulates priming only at long SOAs); participant counts, exact SOAs, item-selection, and statistical results are here.

\subsection{Participants and design}

Two masked-priming experiments, one with English -\mention{ness} nominalizations (49 native speakers: 24 in the short-SOA condition, 25 in the long-SOA condition) and one with Russian -\mention{ost'} nominalizations (61 native speakers: 31 short-SOA, 30 long-SOA) \autocite[1115]{heyer-kornishova-2018-priming-eventually}.

Targets in both experiments were the adjective bases (e.g., \mention{pale, dark}). Primes came in three types:

\begin{itemize}
\tightlist
\item
  the corresponding nominalization as a related prime (e.g., \mention{paleness} preceding \mention{pale});
\item
  a simplex noun as an unrelated control prime;
\item
  the adjective itself as an identity prime.
\end{itemize}

Each experiment used two stimulus-onset asynchronies (SOAs):

\begin{itemize}
\tightlist
\item
  short: 39 ms English, 33 ms Russian;
\item
  long: 77 ms English, 67 ms Russian.
\end{itemize}

\subsection{Items and transparency rating}

Items spanned the full transparency range. At the opaque end, \mention{business} (rated 3.17) is no longer interpretable as \enquote{state of being busy}; at the transparent end, \mention{falseness} (rated 5.75) is fully derivable from \mention{false}; intermediate items occupy the middle of the scale.

Out of 95 -\mention{ness} items pre-tested, 30 were selected for the main experiment to cover the range continuously \autocite[1115]{heyer-kornishova-2018-priming-eventually}. The Russian experiment used a parallel design with 30 -\mention{ost'} nominalizations spanning a similar range.

Transparency was rated by separate groups of native speakers (28 English, 38 Russian) on a 1--7 Likert scale for semantic relatedness.

The two suffixes are highly comparable: both deadjectival, both producing abstract nouns denoting \enquote{state of being X} (e.g., \mention{dark} + \mention{ness} = `darkness'; \mention{bledn-} + \mention{-ost'} = \mention{blednost'} `paleness'), and both productive \autocite[1115]{heyer-kornishova-2018-priming-eventually}.

\subsection{Statistical results}\label{sec:C-3}

At short SOA in both languages, the data shows constant morphological priming across the transparency scale. The interaction between Prime Type and Transparency doesn't reach significance: priming effects are 37 ms in English and 16 ms in Russian, but they aren't modulated by how transparent the prime-target pair is rated \autocite[1117--1119]{heyer-kornishova-2018-priming-eventually}.

At long SOA, priming becomes transparency-modulated: more transparent prime-target pairs (e.g., \mention{paleness-pale}) yield larger priming effects than less transparent pairs (e.g., \mention{business-busy}). The three-way interaction between Prime Type, Transparency, and SOA is significant in English (\mention{t} = 2.10, \mention{p} = 0.036) and marginally so in Russian (\mention{t} = -1.69, \mention{p} = 0.091; the negative sign reflects factor coding rather than a reversed effect, and the direction of transparency-on-priming is consistent across the two languages). The English result establishes the SOA-dependent modulation; the Russian result trends in the same direction without reaching conventional significance \autocite[1117--1119]{heyer-kornishova-2018-priming-eventually}.

\subsection{Rival parallel-processing position}

The serial reading is contested. Feldman et al.'s \autocite*{feldman-etal-2015-must-analysis} \mention{Must analysis of meaning follow analysis of form?} is the canonical parallel-immediate paper: a time-course analysis with a same-target design (each target compared after semantically similar and semantically dissimilar morphologically structured primes within subjects) and SOA treated as continuous from 34 to 100~ms. Their finding is that semantic similarity affects recognition latencies even at the shortest SOAs (34 and 48~ms), with the effect of shared semantics increasing linearly with SOA when long SOAs are intermixed. The serial form-then-meaning ordering predicts no semantic effect at SOAs in the 34--48~ms range; the parallel-immediate reading takes the observed effect as direct evidence against that prediction. Heyer and Kornishova's design (separate item sets at each SOA, transparency as a continuous predictor, two languages) gives the serial reading a different test: the absence of transparency modulation across the item set at the short SOA. The two designs probe different patterns. Same-target wedge (Feldman et al.) tests whether semantic similarity affects latencies on identical targets at very short SOAs; across-item gradient (Heyer and Kornishova) tests whether transparency rating predicts priming magnitude on a continuum. They aren't head-to-head, and the literature is mixed enough (the meta-table on p.~1121 of Heyer and Kornishova surveys 35 prior masked-priming studies) that the present paper takes the serial reading as Heyer and Kornishova's reading, not as a settled disciplinary verdict.
